\section{Methods}
\label{sec:methods}

\subsection{Dataset and cohort}

All experiments use a single curated dataset of B-mode ultrasound images of
schistosomiasis-associated liver fibrosis, comprising 108,709 region-of-interest
crops from 6,373 patients across 35 centers. The study protocol was approved by
the Ethics Review Committee of the Jiangsu Institute of Parasitic Diseases
(Approval ID: JIPD-2019-017), and all participants provided written informed
consent prior to enrolment. Enrolment followed the criteria of the predecessor
cohort study \cite{xu2026sfibai}: participants were recruited from
schistosomiasis-endemic counties in Jiangsu Province, China, covering the full
severity spectrum, including individuals with a confirmed history of
\textit{Schistosoma japonicum} infection, registered advanced schistosomiasis
patients, and non-infected controls. Sex and gender were not collected and were
not considered in the study design.

Each patient is acquired under a standardized six-position anatomical protocol
comprising the subxiphoid sagittal plane (showing the left lobe of the liver),
the left subcostal transverse plane (showing the left lobe and the sagittal
segment of the portal vein), two right subcostal planes (showing the second
porta hepatis with the confluence of three hepatic veins, and the right hepatic
lobe with the right hepatic vein and diaphragm, respectively), and two right
intercostal oblique planes (showing the right hepatic lobe with the portal vein
and gallbladder, and with the right kidney, respectively). One to three images
are acquired per plane, and every image carries a normalized position label
together with a fibrosis score.

Fibrosis severity is annotated on a fine-grained 36-level scale spanning
$0.0$ to $3.5$ in steps of $0.1$. The four clinical strata derive from the
Chinese national grading system for schistosomiasis-associated liver fibrosis
\cite{moh2000schistomanual,moh2006schistocriteria}; their correspondence with
international ultrasonographic guidelines, which were developed primarily for
\textit{S.~mansoni} infection, is described in the WHO-chaired Basel
ultrasonography protocol \cite{richter2025basel}. For clinical reporting the
continuous scale is mapped to the conventional four-grade scale at the
boundaries $0.5$, $1.5$ and $2.5$, with boundary values assigned to the higher
grade. The image label is the maximum annotated score for that image, and the
patient label is the maximum over that patient's images, reflecting the
clinical convention that patient-level severity follows the most advanced
finding.

Annotation was performed in Labelme, with a rectangular bounding box drawn
around each fibrotic region and assigned both the four-grade grouping and the
fine-grained 36-level score. To ensure consistency in applying the 0.1-level
scale, annotating physicians underwent two rounds of unified municipal- and
provincial-level training; the labels were then centrally reviewed by a
coordinating group, with corrective feedback on discrepancies, and a subset of
cases underwent blinded review and secondary evaluation. The target overall
discrepancy rate was below 5\%, and cases with major disagreements were jointly
reviewed to reach a consensus grade.

Quality control of the image data excluded: an acquisition batch collected
under a non-standard position protocol; images whose region-of-interest crop
or annotation box failed validation; duplicate images identified by post-crop
content hashing, including all conflicting duplicates crossing patients,
centers or labels; images without a valid annotation, which were not treated
as negative; and patients with fewer than five of the six standard positions
represented. Splits were generated only after cleaning and deduplication, so
that no patient or duplicate group crosses a split boundary.

The dataset is partitioned at the patient level into 83,722 training images
from 4,906 patients, 20,880 validation images from 1,227 patients, and 4,107
test images from 240 patients. There is no patient overlap across training,
validation and test sets. Centers were shared across splits: all four centers
represented in the test set also appear in the training split. The test set
therefore evaluates generalization to unseen patients within observed centers,
and we refer to it throughout as the patient-held-out test set. Center-balanced
patient-level metrics were used to reduce the influence of differences in center
sample size, rather than to estimate performance on unseen centers.

Table~\ref{tab:cohort} reports the full per-split composition. The test cohort
contains 60 patients in each of F0, F1, F2 and F3, whereas training and
validation follow the cohort distribution, in which F2 is the most frequent
patient-level grade. Because the test cohort was stage-balanced (60 patients per
stage), aggregate performance should not be interpreted as a prevalence-weighted
estimate for a population screening setting. Image-level counts are not exactly
balanced, since patients contribute differing numbers of images.

\subsection{Preprocessing}

All images are resized to $512\times512$ and normalized with ImageNet
statistics. The principal configurations resize by anisotropic stretching; one
control configuration instead uses aspect-preserving letterbox padding, so that
the effect of the resize convention can be separated from the modeling
contributions.

Training augmentation consists of a non-zero multiple-of-$90^\circ$ rotation
applied with probability $0.7$ and a saturation jitter in the range
$0.75$--$1.25$ applied with probability $0.5$. No flipping or cropping is used,
because horizontal flipping is not anatomically meaningful under the six-position
acquisition protocol and cropping would disturb the region-of-interest geometry.
Configurations that are compared pairwise share both sample ordering and
per-sample augmentation draws, so that observed differences do not arise from
divergent data presentation.

For the lesion branch, the weak supervision target is constructed as a union
mask over the annotated boxes in the original region-of-interest coordinate
frame, and is then resized and rotated synchronously with the image using
nearest-neighbour interpolation to preserve label values. Systematic
misalignment between image and mask is treated as a hard failure; single-pixel
boundary discretization differences are tolerated.

\subsection{SynAP-Fib architecture}

\subsubsection{Grading backbone}

The grading pathway follows the SFibAI formulation~\cite{xu2026sfibai}: a
ResNet-50 backbone initialized from \texttt{IMAGENET1K\_V2} weights, followed by
an ordinal head producing a posterior over the 36 discrete fibrosis levels. The
reported continuous prediction is the posterior expectation
\begin{equation}
\hat{y} = \sum_{k=0}^{35} p_k \cdot \frac{k}{10},
\label{eq:score}
\end{equation}
which yields a smooth estimate consistent with the continuous progression of
fibrosis rather than a hard level assignment.

The SFibAI baseline was retrained from scratch under the present
preprocessing, split, and evaluation protocol; consequently, its performance in
this study should not be directly compared with values reported under the
original SFibAI experimental setting.

\subsubsection{Anatomical prior branches}

SynAP-Fib augments this backbone with two anatomical prior branches, shown in
Fig.~\ref{fig:architecture}.

The \emph{position branch} is a six-class head trained with cross-entropy
against the normalized position label. Ground-truth position is used only as a
training supervision signal. At inference the grading pathway reads the
\emph{predicted} position posterior, never the position metadata, so that the
method remains deployable in settings where reliable position metadata is
unavailable.

The \emph{lesion branch} predicts a spatial attention map supervised by the weak
box union target. Only images with a positive fibrosis grade and at least one
annotated box contribute to this objective; F0 images and images lacking boxes
are excluded from the lesion loss.

\subsubsection{Gated residual injection}

Each branch influences grading through the same mechanism. The branch output is
converted into a weighted feature representation and added to the grading
representation as a residual, scaled by a learned gate bounded above by
$\mathrm{gate\_max}=0.5$ and initialized at
$\mathrm{residual\_init\_scale}=0.001$ so that training begins close to the
baseline pathway. The joint configuration superimposes both residuals.

The gradient boundary between the branches and the grading pathway is
bidirectionally detached: gradients from the grading loss do not flow back into
the branch heads through the injection, and gradients from the branch losses do
not propagate into the grading pathway through it. The branches therefore act as
sources of anatomical context rather than as additional gradient paths into the
shared representation. All gate hyperparameters are held fixed across every
configuration, so no configuration is advantaged by gate tuning.

\subsection{Training protocol}

Every configuration is trained for 120 epochs without early stopping, using
AdamW with learning rate $10^{-4}$, weight decay $10^{-4}$, batch size 24, mixed
precision, and a step schedule that multiplies the learning rate by $0.6$ every
15 epochs. All configurations start from the same pretrained initialization,
with modules shared between configurations initialized identically tensor by
tensor under the common seed.

The grading objective is the inherited SFibAI hybrid ordinal loss, which
combines a neighborhood-aware soft-label term, an expectation regression term
and a boundary-sensitive penalty, weighted by $\alpha=1.0$, $\beta=0.02$ and
$\gamma=0.02$ with a soft-label standard deviation of $1.0$. The position and
weak-box objectives each enter the total loss with an outer weight of $0.1$.

All models were trained on a single NVIDIA GeForce RTX 5090 D GPU (CUDA 12.8)
using PyTorch 2.11.0 and torchvision 0.26.0 under Python 3.10, with automatic
mixed precision.

\subsection{Evaluation protocol}

\subsubsection{Clinical objective risk}

Evaluation uses a composite clinical objective risk (COR) that penalizes both
continuous score error and clinically consequential grade errors. Writing
$e = |\hat{y}-y|$ for the absolute score error and
$d = |\mathrm{grade}(\hat{y})-\mathrm{grade}(y)|$ for the grade distance,
\begin{align}
\mathrm{COR} ={}& 0.35\frac{e}{3.5} + 0.15\,\mathbb{I}(e>0.3) + 0.15\,\mathbb{I}(e>0.5) \nonumber\\
             & + 0.25\frac{d}{3} + 0.10\,\mathbb{I}(d\geq2).
\label{eq:cor}
\end{align}
All samples are weighted equally and lower values indicate better performance.

\subsubsection{Aggregation levels}

COR is aggregated at three levels. The image-level risk $R_{\mathrm{image}}$ is
the mean per-image COR. The patient-max risk $R_{\mathrm{patient\text{-}max}}$
takes the maximum over each patient's true and predicted image scores
symmetrically before computing COR, matching the clinical convention that
patient severity follows the most advanced image. The center-balanced risk
$R_{\mathrm{cb}}$ computes patient-max COR within each center and then combines
centers with weights proportional to the square root of each center's patient
count; its purpose is to prevent large centers from disproportionately
dominating the aggregate patient-level evaluation, not to estimate performance
on unseen centers.

The prespecified primary endpoint combines the three levels as
\begin{equation}
R_{\mathrm{final}} = 0.4\,R_{\mathrm{image}} + 0.4\,R_{\mathrm{patient\text{-}max}}
                   + 0.2\,R_{\mathrm{cb}}.
\label{eq:rfinal}
\end{equation}
A patient-median aggregation, computed symmetrically over true and predicted
scores, is retained as a robustness analysis.

\subsubsection{Reported metrics}

Alongside COR we report mean absolute error, accuracy within $\pm0.3$ and
$\pm0.5$, four-grade accuracy and macro-F1. Two error-profile metrics are used.
The excess-error metric beyond a $0.5$ tolerance,
$\mathrm{TMAE} = \mathbb{E}[\max(e-0.5,\,0)]$, is zero for predictions inside
the tolerance band and grows only with the portion of error exceeding it. The
severe error rate is defined on the continuous score as
$\Pr(e>1.0)$. Discrimination and calibration are summarized by macro AUROC,
macro AUPRC, the Brier score and expected calibration error.

For the lesion branch we report Dice and IoU at a threshold of $0.5$ against the
weak box union. These quantify agreement with weak bounding boxes and must not
be read as lesion segmentation ground-truth performance, since no pixel-level
lesion annotation exists in this cohort.

\subsection{Checkpoint selection and reporting discipline}

Checkpoint selection uses the validation set only. Every epoch is fully
evaluated on validation; epochs 1--20 are treated as burn-in and are ineligible
for selection. Within epochs 21--120 the selected checkpoint is the one with the
lowest validation $R_{\mathrm{final}}$, resolving ties by lower validation image
COR and then by earlier epoch. No smoothing, early stopping or selector
resampling is applied. The test set plays no role in checkpoint, threshold or
calibration selection.

The comparison reported here is a single-seed comparison under seed 2026, and we
do not present it as evidence about run-to-run variability. Two metrics are
recomputed post hoc from the frozen per-image predictions and are labelled as
such throughout: stage-wise mean absolute error, and quadratic-weighted Cohen's
$\kappa$ over the 36 levels. Continuous posterior-expectation predictions were
mapped to the nearest one of the 36 valid grades before computing
quadratic-weighted Cohen's $\kappa$; posterior argmax predictions were not used.
